\section{Metodologia}
\label{sec:metodologia}

Esta seção descreve o desenho metodológico do projeto. Os valores numéricos de
hiperparâmetros e suas justificativas são apresentados nas seções de cada
\emph{milestone} e consolidados nos ADRs do repositório.

\subsection{Pré-processamento}
A feature principal é o \emph{retorno logarítmico diário},
$r_t = \ln(P_t / P_{t-1})$, sobre o preço de fechamento ajustado. Adota-se um
\textbf{split temporal antes da normalização}: o \texttt{MinMaxScaler} é ajustado
\emph{apenas} sobre o treino e então aplicado ao teste, evitando vazamento de
informação do futuro. A série é então segmentada em \emph{janelas deslizantes} de
tamanho fixo (ver Seção~\ref{sec:m1}).

\subsection{Modelo}
A arquitetura é um \emph{LSTM-Autoencoder}~\cite{hochreiter1997lstm}: um codificador
LSTM comprime a janela em uma representação latente (gargalo), e um decodificador LSTM
reconstrói a sequência, otimizando o erro quadrático médio (MSE) de reconstrução. Treina-se
\textbf{um modelo por ativo}, condição necessária para a comparação setorial.

\subsection{Detecção}
Uma janela é marcada como anômala quando seu erro de reconstrução ultrapassa um limiar.
Comparam-se duas estratégias:
\begin{itemize}[noitemsep]
  \item \textbf{Estático:} percentil fixo (95) do erro de reconstrução do treino.
  \item \textbf{Dinâmico:} percentil em janela móvel \emph{causal} (apenas o passado),
        sensível a mudanças de regime.
\end{itemize}

\subsection{Avaliação}
Por ser um problema não supervisionado, a avaliação combina duas vias~\cite{liu2025robust}:
\begin{itemize}[noitemsep]
  \item \textbf{Narrativa:} verificação de que as anomalias detectadas correspondem a
        eventos reais conhecidos (e.g., \emph{crash} de março/2020; caso Americanas em
        janeiro/2023).
  \item \textbf{Quantitativa:} injeção controlada de anomalias sintéticas (\emph{price
        shocks} e \emph{volume spikes}) em posições conhecidas, permitindo calcular
        Precision, Recall e F1.
\end{itemize}

\paragraph{Nota metodológica.} O período de treino (2010--2019) contém eventos de forte
volatilidade (recessão de 2014--2016, Operação Lava Jato, \emph{impeachment} de 2016).
A definição de ``normalidade'' é, portanto, relativa, e essa limitação é assumida
explicitamente.
